\section{多层中间表示：让信息在需要时才消失}
\label{sec:mlir}

LLVM IR 适合表达已经具体化的控制流、内存访问和调用约定，但编译器并不总想立刻回答这些问题。对截断点积而言，式~\eqref{eq:clipped-dot} 首先说的是三个结构：逐元素相乘、逐元素限幅、沿索引维归约。若前端立即把它们展开为指针递增和条件跳转，后端仍可能生成好代码，却必须从偶然的循环形状中重新辨认这些结构。多层中间表示（multi-level intermediate representation）的核心思想不是增加更多“必经阶段”，而是延长有用信息的寿命，并在决策真正发生时才将其具体化。MLIR 为这种做法提供了可扩展的操作、类型、属性、接口和转换基础设施\cite{adv-mlir-cgo,adv-mlir-langref}。

% Progressive MLIR lowering / MLIR 渐进式降低
\begin{figure}[htbp]
  \centering
  \begin{tikzpicture}
    \matrix (m) [matrix of nodes,
      nodes={vis/artifact,text width=2.55cm,minimum height=.78cm},
      row sep=2.4mm,column sep=2.5mm,
      column 1/.style={nodes={vis/label,text width=1.05cm,draw=none,fill=none}},
      column 2/.style={nodes={vis/semantic node}},
      column 3/.style={nodes={vis/decision}},
      column 4/.style={nodes={vis/artifact}}] {
      & |[vis/label,font=\bfseries\footnotesize,draw=none,fill=none]| 结构化快照
      & |[vis/label,font=\bfseries\footnotesize,draw=none,fill=none]| 混合快照
      & |[vis/label,font=\bfseries\footnotesize,draw=none,fill=none]| 面向目标的快照 \\
      计算 & {\shortstack{\texttt{linalg} 归约\\截断映射}}
           & {\shortstack{\texttt{vector} 与\\残余 \texttt{linalg}}}
           & {\shortstack{目标内建函数\\LLVM 方言}} \\
      数据 & {\shortstack{\texttt{tensor} 值\\形状属性}}
           & {\shortstack{\texttt{tensor}/\texttt{memref}\\并存}}
           & {\shortstack{缓冲区\\地址空间}} \\
      控制 & 隐式迭代域
           & {\shortstack{\texttt{scf.forall}\\\texttt{scf.for}}}
           & 分支与调用 \\
      目标 & 布局仍抽象
           & {\shortstack{向量宽度\\内存空间}}
           & ABI 与设备操作 \\
    };

    \draw[vis/edge] ([yshift=2mm]m-1-2.north east) --
      node[vis/label,above]{依接口改写} ([yshift=2mm]m-1-3.north west);
    \draw[vis/edge] ([yshift=2mm]m-1-3.north east) --
      node[vis/label,above]{按目标合法化} ([yshift=2mm]m-1-4.north west);

    \draw[vis/semantic,line width=1pt,rounded corners=2pt]
      ([xshift=-2pt,yshift=2pt]m-2-2.north west) rectangle
      ([xshift=2pt,yshift=-2pt]m-5-3.south east);
    \node[vis/label,text=semblue,fill=white,inner sep=1.5pt]
      at ($(m-5-2.south)!0.5!(m-5-3.south)+(0,-.28)$)
      {结构化事实仍可由接口访问};

    \draw[vis/deferred edge] ([yshift=-3mm]m-5-2.south) to[bend right=20]
      node[vis/label,below]{跳过混合态：过早丢失信息}
      ([yshift=-3mm]m-5-4.south);
    \coordinate (x) at ($(m-5-3.south)+(0,-.55)$);
    \draw[vis/cross] ($(x)+(-2mm,-2mm)$) -- ($(x)+(2mm,2mm)$);
    \draw[vis/cross] ($(x)+(-2mm,2mm)$) -- ($(x)+(2mm,-2mm)$);
  \end{tikzpicture}
  \caption{渐进式降低改变的是模块中各种抽象的组成比例：计算、数据、控制和目标约束可以按不同节奏显式化；接口与合法性条件使混合表示仍可组合。}
  \label{fig:mlir-lowering}
\end{figure}


\subsection{混合方言是正常状态}

方言（dialect）是语义命名空间，不是流水线中的独占阶段。一个合法模块可以同时使用 \code{linalg} 描述结构化归约、\code{tensor} 表达值语义、\code{arith} 表达标量比较、\code{scf} 表达循环，并夹杂尚未转换的调用。图~\ref{fig:mlir-lowering} 因而不是一列互斥方框：降低过程改变的是模块中各种抽象的比例，而不是按下开关后把整个程序瞬间换成“下一种 IR”。

对长度 $k$ 的截断点积，可把核心计算写成
\begin{equation}
  p_i=a_i b_i,\qquad
  c_i=\min(12,\max(-8,p_i)),\qquad
  r=\sum_{i=0}^{k-1}c_i .
  \label{eq:mlir-map-reduce}
\end{equation}
前两个式子是同一索引空间上的结构化映射，最后一个式子是以 $0$ 为初值的加法规约。这一描述保存了索引映射、归约维和逐元素区域，因此融合、分块和向量化不必先做循环识别。边界 $0\le k\le8$ 与 $-8\le c_i\le12$ 又给出 $-64\le r\le96$，所以这里把归约重排为加法树不会引入有符号溢出；若删除长度上界，这项工程自由便随之消失。形式化不变量在此不是装饰，而是优化许可。

仅有许多方言并不足以形成系统。接口（interface）让通用变换询问一个操作“拥有什么能力”，而不是枚举它“属于哪个名字”。例如，结构化操作接口暴露迭代器和索引映射，目的地风格接口显式给出输出或初值，内存效果接口说明读写行为，缓冲化接口描述别名和存储改写\cite{adv-mlir-interfaces,adv-mlir-linalg-rationale}。因此，新的领域操作可以接入已有的分块或缓冲化算法；反过来，若一个操作没有声明相应契约，通用算法就不应仅凭外形猜测其语义。这是可扩展性与可推理性之间最重要的平衡。

\subsection{降低是一个合法性问题}

模式改写（pattern rewriting）给出局部等价替换；方言转换（dialect conversion）进一步规定整个转换结束时允许留下什么。令模块中的操作集合为 $O$，转换目标给出合法性谓词（legality predicate）
\begin{equation}
  L_t(o,\Gamma)\in\{\mathsf{true},\mathsf{false}\},
  \qquad
  \operatorname{Legal}_t(M)\iff
  \forall o\in O(M),\ L_t(o,\Gamma_M),
  \label{eq:mlir-legality}
\end{equation}
其中 $t$ 是本次转换目标，$\Gamma_M$ 包含类型、属性和所在区域等上下文。静态合法性只依赖操作种类；动态合法性还可要求函数签名已经转换、布局满足限制或某个属性已经确定。部分转换只要求被声明为非法的操作消失，完全转换则要求范围内每个操作都被目标接纳\cite{adv-mlir-conversion}。于是“中间模块同时含有 \code{linalg}、\code{scf} 和 \code{vector}”并不表示流水线只做了一半；只要式~\eqref{eq:mlir-legality} 对当前目标成立，它就是刻意选择的稳定交接面。

这种设计把一个难以维护的全局转换拆成可组合的责任：规则只需把自己承诺处理的非法操作改写为更接近目标的形式，框架负责反复应用规则并检查终态。代价也很现实。合法集合定义得过宽，错误会被推迟到更难诊断的阶段；定义得过窄，则每个实验性操作都必须立刻获得完整后端支持。好的转换边界不是“越低越严谨”，而是在当前决策所需的信息与下一阶段可接受的复杂度之间取平衡。

\subsection{类型转换需要跨边界的物化}

操作合法化不能脱离类型转换（type conversion）。\code{TypeConverter} 可以把一个源类型映射为零个、一个或多个目标类型；例如，排序 \code{memref} 在 LLVM ABI 边界常展开为分配指针、对齐指针、偏移、各维大小与步长，而不是简单改名为一个 \code{ptr}\cite{adv-mlir-conversion,adv-mlir-llvm-target}。若生产者已转换而消费者仍停留在源方言，两者的 SSA 类型便不再相同。源物化、目标物化与参数物化（materialization）在这个断面构造桥接值，使局部转换仍能保持整个使用--定义链类型正确。

临时的 \code{builtin.unrealized\_conversion\_cast} 可用于推迟桥接细节，但它没有独立的机器语义。允许它永久流入代码生成，相当于把尚未解决的类型义务伪装成普通操作。工程上应把这类 cast 看作“语义债务”：混合 IR 内可以短暂存在，越过最终翻译或 ABI 边界前必须被真正的物化规则消除。这也说明渐进式降低并非放弃严格性；它只是把严格性表述为每个边界各自可检查的合同。

\subsection{缓冲化改变了程序讨论的问题}

张量值回答“这个计算产生什么值”，缓冲区还要回答“该值存在哪里、与谁别名、由谁释放”。缓冲化（bufferization）因此不是换一种类型拼写，而是从值等价进入存储身份。One-Shot Bufferize 沿 SSA 使用--定义链分析读写冲突，尽量将目的地风格操作原地化；当原地更新可能改变仍可观察的旧值时，就必须分配或复制\cite{adv-mlir-bufferization}。

在式~\eqref{eq:mlir-map-reduce} 中，$p_i$ 与 $c_i$ 没有必要物化为完整数组：映射可以同归约融合，每次生成一个 $c_i$ 后立即累加。保持张量和结构化操作直到完成融合，使这个结论直接可见。若过早缓冲化，中间数组的分配、别名与生存期会进入分析，编译器随后还要证明它们可以被删掉。反之，过晚缓冲化也会掩盖片上容量、地址空间和数据搬运等硬约束。最佳位置不是固定的“尽可能晚”，而是：先完成依赖值语义的结构变换，再在存储决策开始影响合法性与成本之前显式化内存。

这就是信息寿命（information lifetime）的工程权衡。高层信息保留越久，通用融合和调度的选择越多；但抽象层越多，接口、规范化规则和转换路径的维护面也越大。低层细节暴露越早，后端越容易落实 ABI 与硬件约束；但许多结构一旦粉碎，就只能通过不可靠的模式识别重建。多层 IR 的价值不在“层数多”，而在让每类事实恰好活到最后一个需要它的决策点。

\subsection{从 MLIR 到 LLVM，再到特定设备}

MLIR 的 \code{llvm} 方言是对应 LLVM IR 语义的一种 MLIR 方言，并不等同于外部 LLVM IR 文件。通常先把模块完全合法化为可翻译的 LLVM 及目标内在方言，再机械地翻译为 LLVM IR；排序 \code{memref} 在这里落实为描述符，调用约定和地址空间也必须具体化\cite{adv-mlir-llvm-dialect,adv-mlir-llvm-target}。两套表示不是竞争关系：MLIR 在领域结构和硬件选择尚未确定时保存信息，LLVM 在这些选择已经具体化后承担成熟的低层优化与代码生成。

异构设备使这条原则更鲜明。AscendNPU IR 在 MLIR 上定义 HFusion、HIVM 与 HACC 等方言：前者面向结构化融合，HIVM 描述 tile、搬运、同步与片上存储，HACC 连接宿主和设备入口\cite{adv-ascend-architecture,adv-ascend-source}。但官方 \code{VecAdd} 示例的起点已经是 HIVM/\code{memref} 层：两个 GM 输入经 \code{hivm.hir.load} 搬入 UB，执行 \code{hivm.hir.vadd}，再由 \code{hivm.hir.store} 写回 GM\cite{adv-ascend-vecadd}。它具体说明了设备边界上的地址空间与数据运动，却不代表任意 \code{linalg} 程序都能自动穿过 HFusion 到达该层。

对截断点积，一个自然的设计是让映射、限幅和归约先保持为结构化操作；随后依据目标向量宽度与 UB 容量形成 tile，将 GM--UB 搬运、局部归约和尾部条件显式化；最后才落实设备入口与宿主调用。这里最难的并非把乘法翻译成某条指令，而是决定部分和放在哪里、怎样合并，以及搬运和同步是否抵消并行收益。官方 \code{VecAdd} 恰好标出一条诚实的知识边界：它展示“低层设备程序怎样写”，没有替我们完成“高层归约怎样调度”。

因此，多层 IR 给出的不是自动优化承诺，而是一种组织决策的办法：结构仍有用时不急于丢弃；存储开始支配可行性时显式缓冲化；越过 ABI 或设备边界前，用合法性谓词消除所有未兑现的抽象。编译器由此不再是一串按名称排列的 Pass，而是一系列有意识的信息降解。
